\documentclass[a4paper, 12pt]{article}
\usepackage[ngerman]{babel}
\usepackage{fontspec}

\title{Bosnien und Herzegowina sowie die Karriere von Edin D\v{z}eko}
\author{Mini-Dokumentation (LaTeX-Abschlussübung)}
\date{November 2026}

\begin{document}

\maketitle
\tableofcontents
\newpage

\section{Einleitung}
Bosnien und Herzegowina ist ein Staat in Südosteuropa auf der Balkanhalbinsel, der durch eine vielfältige Geschichte, beeindruckende Gebirgslandschaften und eine leidenschaftliche Sportkultur geprägt ist. Ziel dieser Mini-Dokumentation ist es, einen kompakten Überblick über das Land zu geben und anschließend die sportliche Laufbahn seines bekanntesten Fußballspielers, Edin D\v{z}eko, anhand von Kennzahlen darzustellen.

\subsection{Bosnien und Herzegowina im Überblick}
Das Land erstreckt sich über eine Fläche von rund $51.197\,\text{km}^2$ und zählt etwa $3,14\,\text{Millionen}$ Einwohner. Die Hauptstadt ist Sarajevo. Politisch gliedert sich der Bundesstaat in die Föderation Bosnien und Herzegowina, die Republika Srpska sowie den Sonderverwaltungsbezirk Br\v{c}ko-Distrikt. Geografisch wird das Landesinnere vom Dinarischen Gebirge dominiert, während im Südwesten bei Neum ein rund $20\,\text{km}$ langer Zugang zum Adriatischen Meer besteht.

\subsection{Edin D\v{z}eko -- Rekordspieler und Stürmerstar}
Edin D\v{z}eko (geboren am 17.~März 1986 in Sarajevo) gilt als der erfolgreichste Fußballspieler in der Geschichte Bosnien und Herzegowinas. Seine Profikarriere begann beim Heimatverein FK \v{Z}eljezni\v{c}ar Sarajevo. Über den tschechischen Klub FK Teplice wechselte er 2007 in die deutsche Bundesliga zum VfL Wolfsburg, wo er 2009 Deutscher Meister und 2010 Torschützenkönig wurde.

Es folgten internationale Top-Stationen bei Manchester City, der AS Rom, Inter Mailand, Fenerbah\c{c}e Istanbul und der AC Florenz, bevor er im Januar 2026 nach Deutschland zurückkehrte und sich dem FC Schalke 04 anschloss. In der bosnisch-herzegowinischen A-Nationalmannschaft absolvierte er zwischen 2007 und seinem Abschiedsspiel im Oktober 2026 insgesamt $152$ Länderspiele und erzielte dabei $73$ Tore, womit er sowohl Rekordspieler als auch Rekordtorschütze seines Landes ist.

\end{document}
